\documentclass[11pt,a4paper]{article}
\usepackage[margin=2.5cm]{geometry}
\usepackage{amsmath,amssymb,amsthm}
\usepackage{booktabs}
\usepackage{microtype}
\usepackage{xcolor}
\definecolor{cA}{HTML}{1F4E8C}   % 010
\definecolor{cB}{HTML}{B3261E}   % 100
\definecolor{cC}{HTML}{1B7A3D}   % 001
\definecolor{cD}{HTML}{B36B00}   % 101
\newcommand{\blk}[2]{\texttt{\textbf{\textcolor{#1}{#2}}}}
% prefix, window, suffix, colour
\newcommand{\slide}[4]{\texttt{\textcolor{black!25}{#1}\textbf{\textcolor{#4}{#2}}\textcolor{black!25}{#3}}}
\usepackage{parskip}
\usepackage[colorlinks=true,linkcolor=blue!50!black,urlcolor=blue!50!black]{hyperref}

\theoremstyle{definition}
\newtheorem{definition}{Definition}
\newtheorem{example}[definition]{Example}
\theoremstyle{plain}
\newtheorem{proposition}[definition]{Proposition}
\newtheorem{theorem}[definition]{Theorem}

\title{Sturmian Sequences}
\author{Robin Langer}
\date{}

\begin{document}
\maketitle
\vspace{-2em}

\section{Windows and splits}

Take an infinite sequence of $0$s and $1$s:
\[
  \ell \;=\; 0\,1\,0\,0\,1\,0\,1\,0\,0\,1\,0\,0\,1\,0\,1\,0 \cdots
\]
Slide a window $m$ letters wide along it and write down every distinct block you
see through it. Let $F(m)$ be how many distinct blocks there are.

Do it. Take $m = 3$ and walk the window along the sixteen letters above:

\begin{center}
\begin{tabular}{llll}
\toprule
position & window & block & \\ \midrule
$0$ & \slide{}{010}{0101001001010}{cA} & \blk{cA}{010} & new \\
$1$ & \slide{0}{100}{101001001010}{cB} & \blk{cB}{100} & new \\
$2$ & \slide{01}{001}{01001001010}{cC} & \blk{cC}{001} & new \\
$3$ & \slide{010}{010}{1001001010}{cA} & \blk{cA}{010} & seen \\
$4$ & \slide{0100}{101}{001001010}{cD} & \blk{cD}{101} & new \\
$5$ & \slide{01001}{010}{01001010}{cA} & \blk{cA}{010} & seen \\
$6$ & \slide{010010}{100}{1001010}{cB} & \blk{cB}{100} & seen \\
$7$ & \slide{0100101}{001}{001010}{cC} & \blk{cC}{001} & seen \\
$8$ & \slide{01001010}{010}{01010}{cA} & \blk{cA}{010} & seen \\
\multicolumn{4}{c}{$\cdots$ and nothing new appears after that $\cdots$} \\
\bottomrule
\end{tabular}
\end{center}

Four distinct blocks, so $F(3) = 4$. Note which of the eight possible
three-letter blocks never appear at all: \texttt{110}, \texttt{011},
\texttt{111} and \texttt{000}. This sequence is far from arbitrary.

Now widen the window by one letter and ask what becomes of each block. Do it:
for each of the four, collect every letter seen following it anywhere in the
sequence.

\begin{center}
\begin{tabular}{llll}
\toprule
block & occurs at & followed by & becomes \\ \midrule
\blk{cA}{010} & $0,3,5,8,11,13$ & $0$ at $0$, \textbf{1} at $3$ & \blk{cA}{0100} \emph{and} \blk{cA}{0101} \\
\blk{cB}{100} & $1,6,9$         & $1$, every time              & \blk{cB}{1001} \\
\blk{cC}{001} & $2,7,10$        & $0$, every time              & \blk{cC}{0010} \\
\blk{cD}{101} & $4,12$          & $0$, every time              & \blk{cD}{1010} \\
\bottomrule
\end{tabular}
\end{center}

Three of the four blocks pass through unchanged. One does
not. \blk{cA}{010} sits at position $0$ followed by a $0$, and at position $3$
followed by a $1$, so it becomes two blocks. Call such a block
\emph{right-special}, and say it \emph{splits}.

The blocks of length four are therefore
\[
  \blk{cC}{0010}, \quad \blk{cA}{0100}, \quad \blk{cA}{0101}, \quad
  \blk{cB}{1001}, \quad \blk{cD}{1010},
\]
five of them, so $F(4) = 5$. We started with four and exactly one split --- and
the two blue ones are the two halves of it.

Every block either passes through or splits in two, so

\begin{definition}
\[
  F(m+1) - F(m) \;=\; \#\{\text{blocks of length } m \text{ that split}\}.
\]
\end{definition}

Here is why the count is worth having. Suppose you are reading $\ell$ one letter
at a time, trying to call the next letter, and you can remember only the last
$m$. Then two occurrences of the same block are indistinguishable to you: you
are in the same situation both times, and whatever rule you are using must say
the same thing at both. If that block splits, the next letter genuinely differs
between the two occurrences, so you are wrong at one of them. Not through bad
judgement --- the letters you can see do not contain the answer.

So counting splits counts the situations in which $m$ letters of memory are not
enough, and a sequence whose splitting stops at some finite $m$ is one you can
predict perfectly from a finite table.

\section{Sturmian sequences}

\begin{definition}
A sequence is \emph{Sturmian} when exactly one block of each length splits:
\[
  F(m) = m + 1 \qquad \text{for every } m \geq 1 .
\]
\end{definition}

The sequence of \S1 is Sturmian: it is the Fibonacci sequence, the fixed point
of $0 \mapsto 01$, $1 \mapsto 0$. We found $F(3) = 4$ and $F(4) = 5$, with
\blk{cA}{010} the block that splits. Here are two more, worked the same way.

\begin{example}[the Pell sequence, $0 \mapsto 001$, $1 \mapsto 0$]
\[
  \ell \;=\; \texttt{001001000100100010010010}\cdots
\]
Its four blocks of length three are \blk{cA}{000}, \blk{cB}{001},
\blk{cC}{010}, \blk{cD}{100}, so $F(3) = 4$. Note $11$ never occurs.

\begin{center}
\begin{tabular}{llll}
\toprule
block & occurs at & followed by & becomes \\ \midrule
\blk{cA}{000} & $6,13,23$   & $1$, every time              & \blk{cA}{0001} \\
\blk{cB}{001} & $0,3,7,10$  & $0$, every time              & \blk{cB}{0010} \\
\blk{cC}{010} & $1,4,8,11$  & $0$, every time              & \blk{cC}{0100} \\
\blk{cD}{100} & $2,5,9,12$  & \textbf{1} at $2$, $0$ at $5$ & \blk{cD}{1001} \emph{and} \blk{cD}{1000} \\
\bottomrule
\end{tabular}
\end{center}

One split, so $F(4) = 5$: \blk{cA}{0001}, \blk{cB}{0010}, \blk{cC}{0100},
\blk{cD}{1000}, \blk{cD}{1001}.
\end{example}

\begin{example}[the rotation of slope $\alpha = 1/\sqrt{2}$]
\[
  \ell \;=\; \texttt{011011011101101110110110}\cdots
\]
Now the $1$s are in the majority --- the density is $0.707$ --- and $00$ never
occurs. The four blocks of length three are \blk{cA}{011}, \blk{cB}{101},
\blk{cC}{110}, \blk{cD}{111}.

\begin{center}
\begin{tabular}{llll}
\toprule
block & occurs at & followed by & becomes \\ \midrule
\blk{cA}{011} & $0,3,6,10$  & $0$ at $0$, \textbf{1} at $6$ & \blk{cA}{0110} \emph{and} \blk{cA}{0111} \\
\blk{cB}{101} & $2,5,9,12$  & $1$, every time               & \blk{cB}{1011} \\
\blk{cC}{110} & $1,4,8,11$  & $1$, every time               & \blk{cC}{1101} \\
\blk{cD}{111} & $7,14,24$   & $0$, every time               & \blk{cD}{1110} \\
\bottomrule
\end{tabular}
\end{center}

Again one split, so $F(4) = 5$: \blk{cA}{0110}, \blk{cA}{0111},
\blk{cB}{1011}, \blk{cC}{1101}, \blk{cD}{1110}.
\end{example}

Three sequences with nothing in common to look at --- one sparse, one balanced,
one dense; two built by substitution and one by rotation --- and in each the
same thing happens. Four blocks of length three, exactly one of them ambiguous,
five blocks of length four. Whatever the slope, there is precisely one situation
at each width in which the letters you can see do not determine the next one.

Those three are not a coincidence, and there is a complete description of the
sequences that behave this way.

\begin{theorem}[Morse--Hedlund]
A sequence $\ell$ is Sturmian if and only if there are an irrational
$\alpha \in (0,1)$ and a $\rho \in [0,1)$ with
\[
  \ell_k \;=\; \lfloor (k+1)\alpha + \rho \rfloor - \lfloor k\alpha + \rho \rfloor
  \qquad \text{for all } k \geq 0 .
\]
The \emph{slope} $\alpha$ is the density of $1$s in $\ell$ and is determined by
$\ell$; the \emph{intercept} $\rho$ only says where in the sequence you started
reading. We do not prove this.
\end{theorem}

So there is one Sturmian sequence for each irrational slope, up to where you
begin. The three worked above are

\begin{center}
\begin{tabular}{llll}
\toprule
sequence & slope $\alpha$ & & \\ \midrule
Fibonacci, $0 \mapsto 01,\ 1 \mapsto 0$ & $1/\varphi^{2} = (3-\sqrt{5})/2$ & $\approx 0.381966$ & no $11$ \\
Pell, $0 \mapsto 001,\ 1 \mapsto 0$     & $1 - 1/\sqrt{2}$                & $\approx 0.292893$ & no $11$ \\
the rotation of \S2                      & $1/\sqrt{2}$                    & $\approx 0.707107$ & no $00$ \\
\bottomrule
\end{tabular}
\end{center}

The last column is not an accident either: a Sturmian sequence contains $11$
exactly when $\alpha > \tfrac12$, and $00$ exactly when $\alpha < \tfrac12$.
Two of these slopes lie below a half and one above.

\section{Why Sturmian sequences make interesting monsters}

Here is a game built from one of these sequences. It is the monster's schedule:
at beat $k$ the monster strikes if $\ell_k = 1$, and stands open if
$\ell_k = 0$.

You never see the current letter before you commit. You see only the letters
already past. Each beat you choose one of two things:

\begin{center}
\begin{tabular}{ll}
\textbf{block}  & nothing happens, whatever the monster does \\[0.3em]
\textbf{attack} & on an open beat, the monster loses $2$ \\
                & into a strike, you lose $5$ \\
\end{tabular}
\end{center}

You start on $20$ points and the monster on $60$. The fight lasts $120$ beats.

Three things follow. You need thirty clean reads to win. You can afford three
mistakes. And blocking forever is a loss, not a draw --- the clock runs out with
the monster untouched.

A player who remembers the last $m$ letters is exactly the thing \S1 describes:
it cannot tell two occurrences of the same block apart, so it plays them
identically, and the blocks that split are precisely the situations where that
costs it. Give such a player a table learned from previous encounters and set it
against monsters of each kind, all with the same density of strikes so that only
the splitting differs:

\begin{center}
\begin{tabular}{lrrrrr}
\toprule
monster & $m=2$ & $m=4$ & $m=6$ & $m=8$ & $m=12$ \\ \midrule
periodic, period $7$ & $0\%$ & $14\%$ & $100\%$ & $100\%$ & $100\%$ \\
Sturmian             & $9\%$ & $47\%$ & $58\%$  & $72\%$  & $86\%$ \\
random               & $0\%$ & $0\%$  & $0\%$   & $0\%$   & $0\%$ \\
\bottomrule
\end{tabular}
\end{center}

\noindent
Win rates over $300$ fights each. Three shapes: a step, a climb, a flat line.

\paragraph{Periodic: nothing left.} Its splitting stops at $m = 6$, and the
table says the same thing. Below six letters the player has almost nothing;
at six it has a complete theory of the monster and never loses again. There is
one threshold, and on the far side of it the game is over permanently --- not
won, but finished, because no further observation can teach anything. A monster
you beat this way you have beaten for good, and the second encounter is the
first one replayed.

\paragraph{Random: nothing to get.} Every block splits at every length, so no
table of any depth is ever a theory of anything. Memory $12$ is worth exactly
what memory $2$ is worth, which is nothing: the player is guessing, guessing
costs $5$ and pays $2$, and it loses every time. Study is not merely
insufficient here, it is pointless, and a player who works that out correctly
concludes there is nothing to work out.

\paragraph{Sturmian: always something, never everything.} Exactly one block of
each length splits. So a player holding $m$ letters has exactly one situation it
cannot call --- and by the nesting of \S2, the ambiguous block at depth $m+1$
sits inside the one at depth $m$. Each extra letter of memory therefore cuts the
outstanding ambiguity into a rarer case without ever removing it. That is
precisely a curve which rises at every depth and reaches $100\%$ at none, and it
is what the middle row measures.

\end{document}